\chapter*{Introduction}
The brain is among the most widely studied subjects of research, and yet many neural behaviours remain elusive to scientists. In recent years, methods adapted from statistical physics, network science, and thermodynamics have been found to explain and provide understanding for computation in the brain. A large part of the brain is devoted to vision, the visual cortex containing between 4-6 billion of the 30 billion neurons of the neocortex, up to 5 times the size of other primates \cite{wandell_visual_2009, andrews_correlated_1997}. Vision is thus an important neural process, and provides one of the first gates of attention to the human brain.

Many approaches to understanding patterns of human gaze have been described. Importantly, there has been success in training large neural networks to predict gaze trajectories \cite{bhandari_modeling_2025, wu_gitnet_2026,abdelrahman_deep_2026,bremer_predicting_2021,ozdel_gaze-guided_2024,akinyelu_convolutional_2022} and saliency maps from an image for human observers \cite{kummerer_modeling_2025,cornia_deep_2016, cornia_predicting_2018, kammerer2016deepgaze2, huang_salicon_2015, bujia_modeling_2022}. Such approaches give often quite accurate results, but they lack the rich understanding of a full model of the network of the brain itself. Theories of biological attention and cognition, such as those of the ``Bayesian brain'', built on active sensing and inference provide a more motivated framework to understand the brain as maintaining a rational causal mental worldview \cite{knill_bayesian_2004, costa_probabilistic_2024, bujia_modeling_2022, tan_theta_2024}. This Bayesian model of the brain has been a fruitful direction to understand the most basic level of neural cognition as sampling through population activity patterns \cite{orban_neural_2016, NIPS2002_a486cd07, beck_probabilistic_2008, zhu_autocorrelated_2024}. Although it has been a debated framework (see \cite{mangalam_myth_2025}), through pioneering work on neural coding \cite{Lange2023-sn, beck_probabilistic_2008}, Karl Friston's introduction of a Free Energy Principle adapted from classical mechanics \cite{Friston2010-aa, }